% GENERATED FILE. Edit canonical lessons, then run npm run generate:books.
% canonical-source-hash: fnv1a64:0105a93577c1e49f
% canonical-lessons: ZH-C252-kaixin, ZH-C252-gandao, ZH-C252-shangxin, ZH-C252-shiwang, ZH-C252-manyi

\chapter{Happy and Sad}
\label{ch:zh-happy-and-sad}
\hlchaptermodality{252}

\begin{chapteropening}
\textbf{By the end of this chapter:} \emph{I can say I am in high spirits, how I feel, and that I am heartbroken, disappointed or satisfied.}

Complete the last lesson of chapter 252: I can say I am in high spirits, how I feel, and that I am heartbroken, disappointed or satisfied.
\end{chapteropening}

\section[\texorpdfstring{\zh{开心} (kāixīn)}{开心 (kāixīn)}]{\zh{开心} (kāixīn) — happy; in high spirits}

\label{lesson:ZH-C252-kaixin}



\begin{quote}
Before the new one: say the Mandarin for a new character, then the Mandarin for to revise a piece of work.
\end{quote}

\subsection*{You'll want to know: \zh{开心}}
\textbf{\zh{开心}} — \emph{kāixīn} — \textquotedblleft{}happy; in high spirits\textquotedblright{}.

\begin{quote}
\textbf{\zh{今天我很开心}} \emph{jīntiān wǒ hěn kāixīn} — I'm in a really good mood today.
\end{quote}

\begin{cousinweb}
\textbf{\zh{开心}} (\emph{kāixīn}) is built from \zh{开} 'to open' and \zh{心} 'the heart', both inherited from Middle Chinese, \emph{khoj} and \emph{sim} (Guangyun, in Baxter's transcription): an open heart.
\end{cousinweb}

\begin{grammarlens}[title={What you've built}]
One more word for describing things.
\end{grammarlens}

\subsection*{Your turn}
\begin{itemize}
\raggedright
  \item \emph{Say it:} \emph{kāixīn}
  \item \emph{Say it:} \emph{kāixīn}, once more
  \item \emph{Say it:} \emph{jīntiān wǒ hěn kāixīn} — I'm in a really good mood today.
  \item \emph{Recall:} say the Mandarin for a new character, then the Mandarin for to revise a piece of work, then say \emph{kāixīn} again
\end{itemize}

\begin{tcolorbox}[breakable,colback=teal!4,colframe=teal!35!black,title={Before you move on}]
What does \zh{开心} mean? (\textquotedblleft{}Happy; in high spirits\textquotedblright{}.) Say it once more.
\end{tcolorbox}
\section[gǎndào]{gǎndào — to feel}

\label{lesson:ZH-C252-gandao}



\begin{quote}
Before the new one: say the Mandarin for to revise a piece of work, then the Mandarin for in high spirits.
\end{quote}

\subsection*{You'll want to know: gǎndào}
\emph{gǎndào} — \textquotedblleft{}to feel\textquotedblright{}.

\begin{quote}
\emph{wǒ gǎndào hěn kāixīn} — I feel really happy.
\end{quote}

Its characters come later; for now it is a word for the ear.

\begin{cousinweb}
\emph{gǎndào} is built from \emph{gǎn} and \zh{到} 'to reach', both inherited from Middle Chinese, \emph{komX} and \emph{tawH} (Guangyun, in Baxter's transcription), which glosses the first 'to move'.
\end{cousinweb}

\begin{grammarlens}[title={What you've built}]
One more word for describing things.
\end{grammarlens}

\subsection*{Your turn}
\begin{itemize}
\raggedright
  \item \emph{Say it:} \emph{gǎndào}
  \item \emph{Say it:} \emph{gǎndào}, once more
  \item \emph{Say it:} \emph{wǒ gǎndào hěn kāixīn} — I feel really happy.
  \item \emph{Recall:} say the Mandarin for to revise a piece of work, then the Mandarin for in high spirits, then say \emph{gǎndào} again
\end{itemize}

\begin{tcolorbox}[breakable,colback=teal!4,colframe=teal!35!black,title={Before you move on}]
What does gǎndào mean? (\textquotedblleft{}To feel\textquotedblright{}.) Say it once more.
\end{tcolorbox}
\section[shāngxīn]{shāngxīn — sad; heartbroken}

\label{lesson:ZH-C252-shangxin}



\begin{quote}
Before the new one: say the Mandarin for in high spirits, then the Mandarin for to feel.
\end{quote}

\subsection*{You'll want to know: shāngxīn}
\emph{shāngxīn} — \textquotedblleft{}sad; heartbroken\textquotedblright{}.

\begin{quote}
\emph{tā hěn shāngxīn} — He's heartbroken.
\end{quote}

Its characters come later; for now it is a word for the ear.

\begin{cousinweb}
\emph{shāngxīn} is built from \emph{shāng} 'to hurt' and \zh{心} 'the heart', both inherited from Middle Chinese, \emph{syang} and \emph{sim} (Guangyun, in Baxter's transcription), which glosses the first 'to harm'.
\end{cousinweb}

\begin{grammarlens}[title={What you've built}]
One more word for describing things.
\end{grammarlens}

\subsection*{Your turn}
\begin{itemize}
\raggedright
  \item \emph{Say it:} \emph{shāngxīn}
  \item \emph{Say it:} \emph{shāngxīn}, once more
  \item \emph{Say it:} \emph{tā hěn shāngxīn} — He's heartbroken.
  \item \emph{Recall:} say the Mandarin for in high spirits, then the Mandarin for to feel, then say \emph{shāngxīn} again
\end{itemize}

\begin{tcolorbox}[breakable,colback=teal!4,colframe=teal!35!black,title={Before you move on}]
What does shāngxīn mean? (\textquotedblleft{}Sad; heartbroken\textquotedblright{}.) Say it once more.
\end{tcolorbox}
\section[shīwàng]{shīwàng — disappointed}

\label{lesson:ZH-C252-shiwang}



\begin{quote}
Before the new one: say the Mandarin for to feel, then the Mandarin for heartbroken.
\end{quote}

\subsection*{You'll want to know: shīwàng}
\emph{shīwàng} — \textquotedblleft{}disappointed\textquotedblright{}.

\begin{quote}
\emph{tā méi lái wǒ hěn shīwàng} — He didn't come, and I'm disappointed.
\end{quote}

Its characters come later; for now it is a word for the ear.

\begin{cousinweb}
\emph{shīwàng} is built from \emph{shī} 'to lose' and \emph{wàng} 'to watch for', both inherited from Middle Chinese, \emph{syit} and \emph{mjangH} (Guangyun, in Baxter's transcription), which glosses the first 'to let slip' and the second 'to watch for someone's return'.
\end{cousinweb}

\begin{grammarlens}[title={What you've built}]
One more word for describing things.
\end{grammarlens}

\subsection*{Your turn}
\begin{itemize}
\raggedright
  \item \emph{Say it:} \emph{shīwàng}
  \item \emph{Say it:} \emph{shīwàng}, once more
  \item \emph{Say it:} \emph{tā méi lái wǒ hěn shīwàng} — He didn't come, and I'm disappointed.
  \item \emph{Recall:} say the Mandarin for to feel, then the Mandarin for heartbroken, then say \emph{shīwàng} again
\end{itemize}

\begin{tcolorbox}[breakable,colback=teal!4,colframe=teal!35!black,title={Before you move on}]
What does shīwàng mean? (\textquotedblleft{}Disappointed\textquotedblright{}.) Say it once more.
\end{tcolorbox}
\section[mǎnyì]{mǎnyì — satisfied; pleased}

\label{lesson:ZH-C252-manyi}



\begin{quote}
Before the new one: say the Mandarin for heartbroken, then the Mandarin for disappointed.
\end{quote}

\subsection*{You'll want to know: mǎnyì}
\emph{mǎnyì} — \textquotedblleft{}satisfied; pleased\textquotedblright{}.

\begin{quote}
\emph{nǐ mǎnyì ma} — Are you happy with it?
\end{quote}

Its characters come later; for now it is a word for the ear.

\begin{cousinweb}
\emph{mǎnyì} is built from \emph{mǎn} 'full' and \emph{yì} 'intent', both inherited from Middle Chinese, \emph{manX} and \emph{'iH} (Guangyun, in Baxter's transcription), which glosses the first 'full' and the second 'will'.
\end{cousinweb}

\begin{grammarlens}[title={What you've built}]
One more word for describing things.
\end{grammarlens}

\subsection*{Your turn}
\begin{itemize}
\raggedright
  \item \emph{Say it:} \emph{mǎnyì}
  \item \emph{Say it:} \emph{mǎnyì}, once more
  \item \emph{Say it:} \emph{nǐ mǎnyì ma} — Are you happy with it?
  \item \emph{Recall:} say the Mandarin for heartbroken, then the Mandarin for disappointed, then say \emph{mǎnyì} again
\end{itemize}

\begin{tcolorbox}[breakable,colback=teal!4,colframe=teal!35!black,title={Before you move on}]
What does mǎnyì mean? (\textquotedblleft{}Satisfied; pleased\textquotedblright{}.) Say it once more.
\end{tcolorbox}
